\section{Preliminaries and notation}\label{sec:pre}

This preliminary section exhibits the most fundamental definitions and notations used in this article.

{\bf About groups.}
Let $G$ be a compact group.
We write $\Irrep(G)$ for the set of equivalence classes of irreducible representations of $G$ and denote by $1 \in \Irrep(G)$ the class of the trivial representation.

{\bf About Hilbert spaces.}
Let $G$ be a compact group.
Furthermore, let $\hH_\sigma$, $\sigma \in \Irrep(G)$, be a~family of Hilbert spaces, and for each $\sigma \in \Irrep(G)$ let $T_\sigma$ be an operators on $\hH_\sigma$.
Throughout this article, we shall freely utilize the fact that $\sigma \mapsto \hH_\sigma$ and $\sigma \mapsto T_\sigma$ can be extended to arbitrary finite-dimensional representations of $G$ by taking direct sums with respect to irreducible subrepresentations.

{\bf About C\Star-dynamical systems.}
Let $\aA$ be a unital C\Star-algebra and let $G$ be a compact group that acts on $\aA$ by \Star-automorphisms $\alpha_g\colon \aA \to \aA$, $g \in G$, such that $G \times \aA \to \aA$, $(g,x) \mapsto \alpha_g(x)$ is continuous.
Throughout this article, we call such data a \emph{C\Star-dynamical system} and denote it briefly by $(\aA,G,\alpha)$.
Moreover, we typically write $\aB := \aA^G$ for the corresponding fixed point algebra.

The conditional expectation $P_1$ onto $\aB$ allows us to define a definite right $\aB$\nobreakdash-valued inner product on $\aA$ by
\begin{gather*}
	\scal{x,y}_\aB := P_1(x^*y) = \int_G \alpha_g(x^*y) \, {\rm d}g,
	\qquad
	x,y \in \aA.
\end{gather*}
The completion of $\aA$ with respect to the induced norm yields a right Hilbert $\aB$\nobreakdash-module, which we denote by $L^2(\aA)$.
The algebra $\aA$ admits a faithful \Star-representation on $L^2(\aA)$ by adjointable operators given by $\lambda\colon \aA \to \End \big( L^2(\aA) \big)$, $\lambda(x)y := x \cdot y$, and consequently we may identify $\aA$ with $\lambda(\aA) \subseteq \End \big( L^2(\aA) \big)$.
Furthermore, for each $g \in G$ we have a unitary operator $U_g$ on $L^2(\aA)$ defined for $x \in \aA \subseteq L^2(\aA)$ by $U_g x := \alpha_g(x)$.
The map $G \ni g \mapsto U_g \in \mathcal{U}\big( L^2(\aA) \big)$ is strongly continuous and implements the \Star-automorphisms $\alpha_g$, $g \in G$, via $\lambda ( \alpha_g(x)) = U_g \lambda(x) U_g^*$ for all $x \in \aA$.
Like every representation of~$G$, the algebra $\aA$ can be decomposed into its isotypic components, let us say, $A(\sigma)$, $\sigma \in \Irrep(G)$, which amounts to saying that their algebraic sum
\begin{gather*}
	\aA_f:= \bigoplus_{\sigma \in \Irrep(G)}^{\rm alg} A(\sigma)
\end{gather*}
is a dense \Star-subalgebra of~$\aA$.
Furthermore, the isotypic components are pairwise orthogonal, right Hilbert $\aB$-submodules of $L^2(\aA)$ such that $L^2(\aA) = \overline{\bigoplus}_{\pi \in \hat G} A(\pi)$.

{\bf About freeness.}
A C\Star-dynamical system $(\aA,G,\alpha)$ is called \emph{free} if the \emph{Ellwood map}
\begin{align*}
	\Phi\colon \ \aA \tensor_{\text{alg}} \aA \rightarrow C(G,\aA),
	\qquad
	\Phi(x\tensor y)(g):=x \alpha_g(y)
\end{align*}
has dense range with respect to the canonical C\Star-norm on $\Cont(G,\aA)$.
This condition was originally introduced for actions of quantum groups on C\Star-algebras by Ellwood~\cite{Ell00} and is known to be equivalent to Rieffel's saturatedness~\cite{Rieffel91} and the Peter--Weyl--Galois condition~\cite{BaCoHa15}.

One of the key tools used in this article is a characterization of freeness that we provided in~\cite[Lemma~3.3]{SchWa17c}, namely that a C\Star-dynamical system $(\aA,G,\alpha)$ is free if and only if for each irreducible representation $(\sigma,V_\sigma)$ of $G$ there is a finite-dimensional Hilbert space~$\hH_\sigma$ and an isometry $s(\sigma) \in \aA \tensor \End(V_\sigma,\hH_\sigma)$ satisfying $\alpha_g(s(\sigma))=s(\sigma) \cdot \sigma_g$ for all $g \in G$.
A rich class of free actions is given by so-called \emph{cleft} actions, which are characterized as follows:
For each irreducible representation $(\sigma,V_\sigma)$ of $G$ there is a unitary element \mbox{$u(\sigma) \in \aA \tensor \End(V_\sigma)$} such that $\alpha_g(u(\sigma))=u(\sigma) \cdot \sigma_g$ for all $g \in G$ (cf.~\cite[Definition~4.1]{SchWa17b}).
